\documentclass[UTF8]{ctexart}
\usepackage{amsmath, amssymb, geometry}
\usepackage{xcolor}
\usepackage{enumitem}
\usepackage{tcolorbox}
\usepackage{fancyhdr}
\usepackage{tikz}
\geometry{a4paper, margin=1in}

% 定义红色环境
\newenvironment{redenv}{\color{red}}{}

% 定义详细解析环境
\newenvironment{detailedanalysis}{%
    \color{red}
    \begin{enumerate}[label=\textbf{第\arabic*步：}, leftmargin=*, align=left, itemsep=1.5ex]
}{%
    \end{enumerate}
}

% 定义概念框
\newenvironment{conceptbox}{%
    \begin{tcolorbox}[colback=red!5, colframe=red!50!black, title=概念解析, fonttitle=\bfseries]
}{%
    \end{tcolorbox}
}

% 定义公式推导环境
\newenvironment{formuladerivation}{%
    \begin{enumerate}[label={\textbf{公式\arabic*：}}, leftmargin=*, align=left]
}{%
    \end{enumerate}
}

% 定义题目和解析的分隔线
\newcommand{\problemrule}{\noindent\rule{\textwidth}{1.5pt}\vspace{-0.8em}}
\newcommand{\analysisrule}{\noindent\rule{\textwidth}{0.8pt}\vspace{-0.8em}}

% 页眉页脚
\pagestyle{fancy}
\fancyhf{}
\fancyhead[C]{九年级上册 · 反比例函数 · 讲义版}
\fancyfoot[C]{\thepage}

\begin{document}

\title{\textbf{\zihao{2} 第一章 反比例函数 深度讲义}}
\author{整理自教材}
\date{}
\maketitle

\section*{一、 反比例函数的概念}

% ========== 第1题 ==========
\problemrule
\section*{【题目1】反比例函数定义辨析}
下列函数是不是反比例函数？若是，请写出它的比例系数。
\begin{enumerate}[label=(\arabic*), itemsep=5pt]
    \item $y = \dfrac{2}{x}$
    \item $y = x^{-1}$
    \item $y = -\dfrac{2}{x}$
    \item $y = \dfrac{1}{2x}$
\end{enumerate}

\begin{redenv}
\analysisrule
\subsection*{逐步解析}

\begin{detailedanalysis}
\item \textbf{问题本质分析}
需要判断给定的函数表达式是否符合反比例函数的定义。
\begin{conceptbox}
\textbf{反比例函数定义}：
形如 $y = \frac{k}{x}$（$k$ 为常数，$k \neq 0$）的函数，叫做反比例函数。
其等价形式有：
1. $y = kx^{-1}$
2. $xy = k$
\end{conceptbox}

\item \textbf{详细求解过程}
\begin{enumerate}[label=(\arabic*)]
    \item \textbf{分析 $y = \dfrac{2}{x}$}：
    符合 $y = \frac{k}{x}$ 的形式，其中 $k=2$。
    \textbf{结论}：是反比例函数，比例系数 $k=2$。

    \item \textbf{分析 $y = x^{-1}$}：
    可以变形为 $y = \frac{1}{x}$，符合 $y = \frac{k}{x}$ 的形式，其中 $k=1$。
    \textbf{结论}：是反比例函数，比例系数 $k=1$。

    \item \textbf{分析 $y = -\dfrac{2}{x}$}：
    可以变形为 $y = \frac{-2}{x}$，符合 $y = \frac{k}{x}$ 的形式，其中 $k=-2$。
    \textbf{结论}：是反比例函数，比例系数 $k=-2$。

    \item \textbf{分析 $y = \dfrac{1}{2x}$}：
    可以变形为 $y = \frac{1/2}{x}$，符合 $y = \frac{k}{x}$ 的形式，其中 $k=\frac{1}{2}$。
    \textbf{结论}：是反比例函数，比例系数 $k=\frac{1}{2}$。
\end{enumerate}

\item \textbf{易错点全面解析}
\begin{itemize}
    \item 对于 $y=\frac{1}{2x}$，容易误判为 $k=1$，忽略了分母中的系数 2 实际上属于比例系数的一部分 $k=1/2$。
    \item 对于 $xy=k$ 的形式，要明确 $k$ 必须是不为 0 的常数。
\end{itemize}
\end{detailedanalysis}
\end{redenv}

% ========== 第2题 ==========
\problemrule
\section*{【题目2】应用建模：进水速度与时间}
已知某空游泳池的容积为 $270 \text{ m}^3$，用恰当的函数表达式来表示进水速度 $v (\text{m}^3/\text{h})$ 与注满该游泳池所需时间 $t (\text{h})$ 之间的关系。

\begin{redenv}
\analysisrule
\subsection*{逐步解析}

\begin{detailedanalysis}
\item \textbf{问题本质分析}
本题考查反比例函数在实际问题中的应用。
基本数量关系：工作总量 = 工作效率 $\times$ 工作时间。
在此题中，容积（总量）= 进水速度 $\times$ 时间。

\item \textbf{详细求解过程}
\begin{enumerate}[label=\alph*.]
    \item \textbf{建立等量关系}：
    由题意知：$v \cdot t = 270$。
    \item \textbf{变形为函数表达形式}：
    因为要求表示 $v$ 与 $t$ 的关系（通常 $v$ 随 $t$ 变化，或视具体语境而定，此处一般写成 $v$ 关于 $t$ 的函数，或 $t$ 关于 $v$ 的函数均可，通常写成 $v = \dots$）：
    \[ v = \frac{270}{t} \]
    \item \textbf{确定自变量范围}：
    在实际问题中，时间 $t$ 必须大于 0。
\end{enumerate}

\item \textbf{最终答案}
\[ v = \frac{270}{t} \quad (t > 0) \]
\end{detailedanalysis}
\end{redenv}

% ========== 第3题 ==========
\problemrule
\section*{【题目3】反比例函数的性质计算}
已知反比例函数 $y = -\dfrac{6}{x}$。
\begin{enumerate}[label=(\arabic*), itemsep=5pt]
    \item 写出这个函数的比例系数和自变量的取值范围；
    \item 求当 $x = -3$ 时函数的值；
    \item 求当 $y = -2$ 时自变量 $x$ 的值。
\end{enumerate}

\begin{redenv}
\analysisrule
\subsection*{逐步解析}

\begin{detailedanalysis}
\item \textbf{详细求解过程}
\begin{enumerate}[label=(\arabic*)]
    \item \textbf{识别参数}：
    对比 $y = \frac{k}{x}$，可知 $k=-6$。
    分母 $x$ 不能为 0，所以自变量取值范围是 $x \neq 0$。
    
    \item \textbf{代入求值}：
    当 $x = -3$ 时，
    \[ y = -\frac{6}{-3} = 2 \]
    
    \item \textbf{解方程}：
    当 $y = -2$ 时，
    \[ -2 = -\frac{6}{x} \]
    两边同时乘以 $x$ (且 $x \neq 0$)：
    \[ -2x = -6 \]
    \[ x = 3 \]
\end{enumerate}

\item \textbf{归纳总结}
已知 $x$ 求 $y$ 或已知 $y$ 求 $x$，都是代入公式进行简单的代数运算。注意符号的处理。
\end{detailedanalysis}
\end{redenv}

\newpage
\section*{二、 反比例函数的图象与性质}

% ========== 第4题 ==========
\problemrule
\section*{【题目4】图象上的点与性质}
已知点 $(3, -3)$ 在反比例函数 $y = \dfrac{k}{x}$ 的图象上。
\begin{enumerate}[label=(\arabic*), itemsep=5pt]
    \item 求这个函数的表达式；
    \item 判断点 $A(-1, 9)$，$B(-3, 2)$ 是否在这个函数的图象上；
    \item 这个函数的图象位于哪些象限？函数值 $y$ 随自变量 $x$ 的增大如何变化？
\end{enumerate}

\begin{redenv}
\analysisrule
\subsection*{逐步解析}

\begin{detailedanalysis}
\item \textbf{解题策略构建}
1. \textbf{待定系数法}：利用已知点坐标求出 $k$。
2. \textbf{点在图象上的判定}：验证点的坐标是否满足函数解析式（即 $x \cdot y$ 是否等于 $k$）。
3. \textbf{性质分析}：根据 $k$ 的正负判断图象位置和增减性。

\item \textbf{详细求解过程}
\begin{enumerate}[label=(\arabic*)]
    \item \textbf{求解析式}：
    将点 $(3, -3)$ 代入 $y = \frac{k}{x}$：
    \[ -3 = \frac{k}{3} \Rightarrow k = 3 \times (-3) = -9 \]
    所以函数表达式为 $\boldsymbol{y = -\dfrac{9}{x}}$。

    \item \textbf{判定点 $A, B$}：
    \begin{itemize}
        \item 对于 $A(-1, 9)$：$x \cdot y = (-1) \times 9 = -9$。等于 $k$，所以 \textbf{点 A 在图象上}。
        \item 对于 $B(-3, 2)$：$x \cdot y = (-3) \times 2 = -6$。不等于 $k (-9)$，所以 \textbf{点 B 不在图象上}。
    \end{itemize}

    \item \textbf{性质判定}：
    因为 $k = -9 < 0$：
    \begin{itemize}
        \item \textbf{位置}：图象位于\textbf{第二、四象限}。
        \item \textbf{增减性}：在\textbf{每一个象限内}，函数值 $y$ 随 $x$ 的增大而\textbf{增大}。
    \end{itemize}
\end{enumerate}

\begin{conceptbox}
\textbf{反比例函数 $y=\frac{k}{x}$ 的性质}：
\begin{itemize}
    \item 当 $k>0$ 时：图象在一、三象限；在每个象限内 $y$ 随 $x$ 增大而\textbf{减小}。
    \item 当 $k<0$ 时：图象在二、四象限；在每个象限内 $y$ 随 $x$ 增大而\textbf{增大}。
\end{itemize}
\textbf{注意}：描述增减性时，必须强调“在每一个象限内”，不能说“在定义域内”。
\end{conceptbox}

\end{detailedanalysis}
\end{redenv}

% ========== 第5题 ==========
\problemrule
\section*{【题目5】交点问题与参数求解}
正比例函数 $y=x$ 的图象与反比例函数 $y = \dfrac{k}{x}$ ($k$ 为常数，$k \neq 0$) 的图象的一个交点的横坐标是 2。求当 $x = -3$ 时，反比例函数 $y = \dfrac{k}{x}$ 的对应函数值。

\begin{redenv}
\analysisrule
\subsection*{逐步解析}

\begin{detailedanalysis}
\item \textbf{详细求解过程}
\begin{enumerate}[label=\textbf{Step \arabic*.}]
    \item \textbf{求交点坐标}：
    交点在正比例函数 $y=x$ 上，且横坐标 $x=2$。
    代入得 $y=2$。
    所以交点坐标为 $(2, 2)$。

    \item \textbf{求 $k$ 值}：
    交点 $(2, 2)$ 也在反比例函数 $y = \frac{k}{x}$ 上。
    \[ k = x \cdot y = 2 \times 2 = 4 \]
    所以反比例函数为 $y = \dfrac{4}{x}$。

    \item \textbf{求函数值}：
    当 $x = -3$ 时，代入解析式：
    \[ y = \frac{4}{-3} = -\frac{4}{3} \]
\end{enumerate}

\item \textbf{最终答案}
当 $x = -3$ 时，对应的函数值为 $-\dfrac{4}{3}$。
\end{detailedanalysis}
\end{redenv}

\end{document}
